% This is free and unencumbered software released into the public domain.
% 
% Anyone is free to copy, modify, publish, use, compile, sell, or
% distribute this software, either in source code form or as a compiled
% binary, for any purpose, commercial or non-commercial, and by any
% means.
% 
% For more information, please refer to <http://unlicense.org/>



% Permet de formater la liste des examinateurs
\newcommand{\BXexam}[5]{#1 \textsc{#2} \dotfill & #3, #4 \dotfill & #5 \\}

% Produit la page de garde et définit le style du reste du document (en-têtes,…)
\newcommand{\BXcouverture}{
	\pagestyle{empty}
	\newgeometry{marginratio=1:1,vmarginratio=1:1,top=4.0cm,bottom=0.5cm}
	\begin{titlepage}
		\begin{center}
			% Pour un meilleur rendu, commenter les deux lignes suivantes. :p
			\vspace*{-96pt}
			\logoLaBRI \hfill \logoUB \hspace*{-32pt} \\
			\vfill
			\Huge  {\bf THÈSE}
			\vfill
			\large PRÉSENTÉE À
			\vfill
			\huge  {\bf L'UNIVERSITÉ DE BORDEAUX}
			\vfill
			\Large ÉCOLE DOCTORALE {\BXed}
			\vfill
			\Large par {\bf {\BXnom}}
			\vfill
			\Large POUR OBTENIR LE GRADE DE
			\vfill
			\LARGE {\bf DOCTEUR}
			\vfill
			\large SPÉCIALITÉ : {\BXspe}
			\vfill
			\rule{4em}{1pt}
			\vfill
			\Large {\bf {\BXsujet}}
			\vfill
			\rule{4em}{1pt}
			\vfill
			Sous la direction de : \Large {\bf {\BXdirector}}
			\vfill
			\rule{4em}{1pt}
			\vfill
		\end{center}
		\vfill
		\noindent {\bf Date de soutenance :} \quad \BXdate
		\vfill
		\noindent {\bf Devant la commission d'examen composée de :} \\
		\begin{tabular}{lll@{}}
			\BXexamlist
		\end{tabular}
		\vfill
		\centering
		\large \rule[0.2em]{0.5em}{1pt} ~{\BXannee}~ \rule[0.2em]{0.5em}{1pt}
                \vfill
                \centering
                \logoRhoban
	\end{titlepage}
	\restoregeometry
	\cleardoublepage
	
	\pagestyle{fancy}
	
	\renewcommand{\footrulewidth}{0.4pt}
	\renewcommand{\chaptermark}[1]{\markboth{\thechapter.\ ##1}{}}
	\renewcommand{\sectionmark}[1]{\markright{\thesection.\ ##1}}
	
	\fancyhead[LE]{}
	\fancyhead[RE]{\slshape \rightmark}
	\fancyhead[LO]{\slshape \leftmark}
	\fancyhead[RO]{}
	
	\fancyfoot[C]{}
	\fancyfoot[LE]{\thepage}
	\fancyfoot[RE]{\BXnomcourt}
	\fancyfoot[LO]{\BXsujetcourt}
	\fancyfoot[RO]{\thepage}
	
	\paragraph*{Résumé} \BXresume
	
	\paragraph*{Title} \BXtitle
	
	\paragraph*{Abstract} \BXabstract
	
	\paragraph*{Keywords} \BXkeywords
	
	\paragraph*{Mots-clés} \BXmotscles
	
	\paragraph*{Laboratoire d'accueil} \BXlabo
}

% Logo UB
\newcommand{\logoUB}{\includegraphics[height=74pt]{../template/logo_ubx.jpg}}

% Logo Rhoban
\newcommand{\logoRhoban}{\includegraphics[height=4cm]{../media/rhoban_logo.png}}

%NOTE : Le strict respect de la charte graphique de l'Université de Bordeaux
% imposerait la présence d'un trait oblique bleu à 30° coupant le logo de l'UB
% (non inclus dans l'image qui m'a été fournie). Il supposerait également que
% le logo du laboratoire soit d'une couleur marron uniforme, l'emploi exclusif
% de la police Times New Roman 11pt ainsi que diverses autres consignes tout
% aussi ubuesques et disponibles à l'URI suivante :
%https://personnels.u-bordeaux.fr/Pratique/Communiquer/Logo-et-charte-graphique2
%
% Aussi, ce modèle de document ne respecte quasiment aucune des spécifications
% de la Direction de la Communication de la Nouvelle Université de Bordeaux.
%
% WONTFIX - « It's not a bug. It's a feature. »

% Logo LaBRI
\newcommand{\logoLaBRI}{
	\begin{tikzpicture}[scale=.06]
	\pgfmathsetmacro\bleulabriR{0/255}
	\pgfmathsetmacro\bleulabriG{114/255}
	\pgfmathsetmacro\bleulabriB{185/255}
	
	\pgfmathsetmacro\vertlabriR{10/255}
	\pgfmathsetmacro\vertlabriG{177/255}
	\pgfmathsetmacro\vertlabriB{134/255}

	\pgfmathsetmacro\rougelabritopR{162/255}
	\pgfmathsetmacro\rougelabritopG{47/255}
	\pgfmathsetmacro\rougelabritopB{50/255}
	
	\pgfmathsetmacro\rougelabrimiddleR{178/255}
	\pgfmathsetmacro\rougelabrimiddleG{117/255}
	\pgfmathsetmacro\rougelabrimiddleB{99/255}
	
	\pgfmathsetmacro\rougelabribottomR{230/255}
	\pgfmathsetmacro\rougelabribottomG{206/255}
	\pgfmathsetmacro\rougelabribottomB{196/255}
	
	
	\definecolor{bleulabri}{rgb}{\bleulabriR, \bleulabriG, \bleulabriB}
	\definecolor{vertlabri}{rgb}{\vertlabriR, \vertlabriG, \vertlabriB}
	\definecolor{rougetoplabri}{rgb}{\rougelabritopR, \rougelabritopG,
	\rougelabritopB}
	\definecolor{rougemiddlelabri}{rgb}{\rougelabrimiddleR,
	\rougelabrimiddleG, \rougelabrimiddleB}
	\definecolor{rougebottomlabri}{rgb}{\rougelabribottomR,
	\rougelabribottomG, \rougelabribottomB}
	
	\def\rectanglepath{-- ++(0cm,38.5cm)-- ++(10cm,0cm)--
	++(0cm,-30.3cm)-- ++(13cm,0cm)-- ++(-2cm,-8.2cm)-- cycle}
	
	
	\shade[top color=rougetoplabri!95!black,bottom
	color=rougebottomlabri!15!white!95!black,middle
	color=rougemiddlelabri!90!white!98!black,shading=axis,
	shading angle=-135] (19.25, 19.25) circle (6.75cm);
	
	\fill[fill=bleulabri] (0,0) \rectanglepath;
	\fill[rotate=180,fill=vertlabri] (-38.5,-38.5) \rectanglepath;
	\end{tikzpicture}
}

% Commandes à redéfinir :
%% Page de garde
\newcommand{\BXed}{}
\newcommand{\BXspe}{}
\newcommand{\BXnom}{}
\newcommand{\BXsujet}{}
\newcommand{\BXdirector}{}
\newcommand{\BXannee}{}
\newcommand{\BXdate}{}
\newcommand{\BXexamlist}{}
%% Pages liminaires
\newcommand{\BXresume}{}
\newcommand{\BXtitle}{}
\newcommand{\BXabstract}{}
\newcommand{\BXmotscles}{}
\newcommand{\BXkeywords}{}
\newcommand{\BXlabo}{}
%% En-têtes
\newcommand{\BXnomcourt}{}
\newcommand{\BXsujetcourt}{}
